\section{设计规范}
\label{sec:principles}

\sysname{} 的开发由一组书面设计规范约束。规范的价值不在于表述本身，而在于每一条都对应一个能由程序检查的不变量、一个执行机制和至少一个测试；违反规范的改动会在测试或流水线闸门上失败。表 \ref{tab:principles} 汇总了九条核心规范，本节其余部分解释其中几条背后的取舍，以及流水线层面另外五条决策。

\begin{table}[t]
\centering
\caption{核心设计规范、对应的不变量与执行机制。}
\label{tab:principles}
\small
\begin{tabularx}{\linewidth}{@{}p{2.4cm}X p{2.9cm}@{}}
\toprule
规范 & 不变量 & 执行机制 \\
\midrule
P1 LLM 不做算术 & 文稿中的任何派生数（增速、占比、约数）都必须出现在本次运行的计算产物中 & 数字护栏 \code{guard} \\
P2 每个数字可追溯 & 抽取事实带 $\sigma$；模型抄的数字必须按数字边界出现在所写页码的原文中；外部数字的摘句必须包含该数字 & \code{find\_evidence}、\code{verify-facts}、\code{verify-sources} \\
P3 双通道不静默取舍 & 规则通道与模型通道取值不一致时生成问题记录，不得只保留一方而不留痕 & \code{merge()}；抽取闸门 \\
P4 一张真值表 & 估值与纠错只读 $T$；抽取存在拒收、冲突、勾稽失败或核心缺失时，流水线在估值前停止 & 流水线阻断 \\
P5 全程追踪 & 每次文件读取与产物写出记录 SHA-256；每步、每次工具调用与模型调用写入事件流 & \code{RunTracer} \\
P6 单位约定 & 事实用元、估值用百万元、比率用百分数值 & 引擎与护栏统一换算 \\
P7 误报控制优先 & 纠错容差取展示精度的半个末位；错误类型必须找到精确的替代解释才下结论；预测与缺真值断言判“无法核实” & \code{verify\_claim} \\
P8 事实→推论→观点 & 报告按“披露事实 / 程序推论 / 作者观点”分层，估值写明假设、适用范围与风险 & 写作技能；出处分类 \\
P9 LLM 输出加护栏 & 任何让模型改写报告的功能都必须加数字锁 & \code{numlock} \\
\bottomrule
\end{tabularx}
\end{table}

\subsection{为什么连约数也不许心算}
P1 最直观的违反是模型写错增速，但更常见、也更难发现的是“约三倍”“超过四成”这类约数。它们看起来无害，读者却会把它们当作事实引用。我们的规则是：如果约数有价值，就让程序算出精确值并写入产物，文稿引用产物；如果没有价值，就删掉。这条规则让护栏可以对文稿中的每个数字一视同仁，不需要区分“精确数”与“示意数”。

\subsection{为什么双通道冲突要交人工}
抽取由模型按页读原文完成，规则通道独立取数做交叉校验。这样的设计面临一个诱惑：冲突时自动采信“更可信”的一方。我们拒绝这样做，因为冲突本身就是最有价值的信号——它往往意味着合并报表与母公司报表混用、续表表头错配、单位声明识别错误。P3 要求冲突以问题记录的形式进入抽取闸门，阻断后续估值，直到有人（或模型回到原文）解释清楚。

\subsection{为什么误报控制优先于召回}
研报纠错的误报代价远高于漏报：一个误报会让使用者怀疑整份核查结果，并且需要人工逐条排除。P7 因此规定：容差按研报展示精度计算，例如“58.62 亿元”的容差是 $\pm 0.005$ 亿元，四舍五入不算错；错误类型（单位、期间、口径、方向）只有在找到一个精确匹配的替代解释时才能下结论，否则一律归为“数值错误”；预测性表述与缺少真值的断言判为“无法核实”，不计为错误。

\subsection{流水线层面的五条决策}
\label{sec:pipeline-decisions}
上述规范面向计算引擎。研报生成流水线在此基础上另有五条决策，它们来自对早期产物的人工质检：

\begin{designbox}{D1 先搜集、后分析}
早期研报中，行业与竞争章节大量出现模型凭记忆写出的数字。我们把外部资料搜集提前为第一个阶段，由独立的研究助理会话完成并登记来源；后续会话只能引用已登记的外部事实。
\end{designbox}
\vspace{0.4ex}
\begin{designbox}{D2 上下文隔离、只靠文件交接}
写作、复核、评审若共享对话，复核者会继承写作者的思路，评审者会看到写作过程而变得宽容。每个角色是一次独立的无头会话，只读写工作区文件；工作区建在仓库目录树之外，看不到其他公司的产物。
\end{designbox}
\vspace{0.4ex}
\begin{designbox}{D3 评分标准对生成方不可见}
评分标准一旦进入生成方的上下文，系统就会朝评分条目优化，而不是朝更好的研报优化。评分标准与完整评测结果保存在工作区之外，回传给生成方的问题单只含理由与证据，不含条目原文与分值。
\end{designbox}
\vspace{0.4ex}
\begin{designbox}{D4 不预设行业框架}
技能说明只给方法与约束（例如“核心驱动至少两环定量推导”），不给“白酒看吨价、电池看 GWh”这类行业模板。模型读完年报与资料包后自行推导该公司的收入与利润驱动。预设模板会让不同公司的报告长得一样，也会让模型在模板不适用时硬套。
\end{designbox}
\vspace{0.4ex}
\begin{designbox}{D5 文风不做程序拦截}
我们曾考虑用正则规则拦截“不是……而是……”“相互印证”等机器腔句式，最终放弃：同样的意思换个说法照样是机器腔，规则只会催生规避。文风交给一个没有参与写作的复核编辑会话改写，数字锁保证它只能动措辞和结构。
\end{designbox}
